\section{The structure shared by the examples}
\label{sec:structure}
The examples use one obstruction in two different ways. For the determined
graphs, it proves maximality by computing the entire monotone polar. For
the selected graphs, it first excludes an entire primal region from the
polar of a countable seed; the extension argument then supplies a maximal
graph with a lower pairing bound. The latter argument adds a choice of
extension, not a new kind of endpoint obstruction.

\subsection{Actual labels and the completed profile}
The coordinates $r(a)$ and $ma$ preserve the symmetric part of the graph
pairing:
\[
\pair{La-Lb}{a-b}=|r(a)-r(b)|^2-\|ma-mb\|_2^2.
\]
Consequently, a Lipschitz curve in the $(r,m)$ coordinates supplies
monotonicity. The complete kernel fibres supply different information:
they constrain every polar query to the same scalar comparison, rather
than merely proving inequalities between the chosen parameter rows.
Both ingredients are necessary in the proofs above. A parameter curve
alone is not an assertion of maximality for its image in $X\times X^*$.

In the architecture of Theorem~\ref{thm:endpoint-criterion}, the precise
criterion is Eq.~\eqref{eq:abs-endpoint-exclusion}. The comparison space
$\R\oplus H_0$ contains the middle segment $(t,\eta)$, $|t|\le1$,
but a point on that segment gives a graph point only if an actual
$p\in X^*$ realizes both $r(p)=t$ and $Mp=(t,\eta)$. The theorem proves
this statement in both directions. In the concrete construction,
$mp\in\ell^1$ rules out that realization. Thus passing from actual
block sums to their $\ell^2$ limits preserves the comparison geometry
but not the existence of dual labels. No assertion that $\ell^2$ is
the continuous dual of the original primal space is involved.

The following consequence makes the loss of dual control visible in
the original Banach spaces, without relying on a picture in the
comparison space.

\begin{proposition}[Convergent primals and unbounded dual labels]
\label{prop:boundary-sequences}
There are sequences of rows common to $A$ and every selected
$\mathcal A$ of Theorem~\ref{thm:nonconstructive} whose primal
coordinates converge in $c_0(I)$ to points outside both domains,
while the norms of their dual coordinates tend to infinity.
There are also such sequences common to $T$ and the corresponding
$\mathcal T$ on standard $\ell^1$, with convergence in its norm
and divergence in the full $\ell^\infty$ dual norm.
\end{proposition}
\begin{proof}
Use $a_n^\pm=a^{\pm(1+1/n)}$ and $x_n^\pm=La_n^\pm$. These rows
belong to $S_0$, hence to both first graphs on $c_0$.
Writing $t=\pm(1+1/n)$ and $\epsilon=\sgn t$, direct substitution in
Eqs.~\eqref{eq:D4}--\eqref{eq:D9} gives the block-zero coordinates
\[
(La^t)_{0,1}=-2t-2\epsilon,\qquad
(La^t)_{0,2}=-3t-2\epsilon,\qquad
(La^t)_{0,3}=-t-\epsilon,
\]
with every later coordinate in that block zero. In each block $b\ge1$,
only the first primal coordinate is nonzero, and it is
$-\omega_{1/n}(b)$. The convergence in $\ell^2$ of these tails implies
convergence in the supremum norm. The two limits $x_\pm$ therefore have
block-zero coordinates $\mp(4,5,2)$ and positive-block first coordinates
$-1/(4b)$. They belong to global $c_0(I)$, and
$\pair{x_\pm}{g}=\pm1$. Neither limit belongs to $\dom A$ or
$\dom\mathcal A$: the proofs of Theorems~\ref{thm:C} and
\ref{thm:nonconstructive} exclude every primal with
$|\pair{x}{g}|\le1$.

On the other hand,
\[
\|a_n^\pm\|_1=2(1+1/n)+1+\|\omega_{1/n}\|_1\longrightarrow\infty.
\]
To verify the divergence, fix any $N$. The sum of the first $N$ tail
coordinates tends to $\sum_{b=1}^N1/(4b)$. These partial harmonic sums
are unbounded as $N\to\infty$, so the full nonnegative tail sums tend
to infinity. This argument does not exchange two infinite sums.

For each sign choose a lift $u_\pm$ of $x_\pm$ through $Q$.
The norm-$2$ lifting bound gives $v_n^\pm$ with
$Qv_n^\pm=x_n^\pm-x_\pm$ and
$\|v_n^\pm\|_1\le2\|x_n^\pm-x_\pm\|_\infty$.
Then $u_n^\pm=u_\pm+v_n^\pm\to u_\pm$ in $\ell^1$, and
$(u_n^\pm,Q^*a_n^\pm)$ belongs to both pullback graphs. The limits
are outside both domains because their $Q$ images are $x_\pm$.
Equation~\eqref{eq:L3} gives
$\|Q^*a_n^\pm\|_\infty=\|a_n^\pm\|_1\to\infty$.
Only controlled individual lifts were used, not a continuous or linear
right inverse.
\end{proof}

\subsection{A lower pairing bound and different second factors}
The examples separate boundedness of the dual labels from control of
their pairing with their own primals. Proposition~\ref{prop:boundary-sequences}
shows that dual norms can diverge along convergent primals. Nevertheless,
the determined graph has the exact identity
\[
\pair{La}{a}=r(a)^2-1-W(|r(a)|-1)>-\sigma,
\]
and the selected graph retains the lower bound $-\sigma$ on every row.
Together with $\|x\|>1/2$, this proves the finite radial bound; it is
not a premise inferred from a conjectured failure of a sum theorem.

The second factors use that control in different ways. A rank-one factor
adds $\lambda r(a)^2$ on the determined graph. For $\lambda\ge\sigma$
the origin is a polar point of the sum. For $0<\lambda<\sigma$, the
finite-support query in Eq.~\eqref{eq:parameter-query} replaces it;
the full computation in Proposition~\ref{prop:all-strengths} is needed
at those strengths. A norm subdifferential instead adds exactly
$\lambda\|x\|$ for every one of its values, so that the lower bound
and the missing primal alone suffice. This is why it also works for
the unspecified first factor in Theorem~\ref{thm:nonconstructive}.
\ifdefined\expandedversion
The bounded-domain normal-cone example in Corollary~\ref{cor:normal}
uses the closed convex set $C_\rho$, which contains $0$. Testing its
normal inequality at $0$ gives $\pair{x}{n}\ge0$ for every
$x\in C_\rho$ and every $n\in N_{C_\rho}(x)$. The explicit interior
witness there verifies the ordered domain condition. That example is a change
of second factor, not a new independent construction of the first.
\fi

The extension proof identifies the additional requirement behind the
nonconstructive family. A lower bound on a seed is not automatically a
lower bound on its full polar, and inclusion-maximality among graphs
obeying that bound is not automatically global maximality. The scaled
polar point in Lemma~\ref{lem:capped-extension} closes precisely this
gap. At the same time, the seed excludes all dual values above the
forbidden primal slab, so every resulting maximal extension still omits
zero. The added row in Eq.~\eqref{eq:nonconstructive-seed} forces the
new factor to differ from the original one while preserving both
properties.

These distinctions delimit the representation result as well as its
use. Theorem~\ref{thm:endpoint-criterion} characterizes maximality for
the stated full-fibre architecture, not for arbitrary monotone operators.
The quotient construction changes the Banach carrier and preserves the
full-dual certificate, not the underlying mechanism. The parameter and
partner variations describe consequences of that mechanism. None of
these statements identifies independent discoveries merely by counting
the resulting counterexamples.
